\documentclass[11pt, a4paper]{article}
\usepackage[a4paper, margin=2.5cm]{geometry}
\usepackage{fontspec}
\usepackage[english, bidi=basic, provide=*]{babel}
\babelprovide[import, onchar=ids fonts]{english}
\babelfont{rm}{Noto Sans}
\usepackage{amsmath}
\usepackage{booktabs}
\usepackage{xcolor}
\usepackage{hyperref}

\definecolor{accent}{HTML}{1E3A8A}
\definecolor{muted}{HTML}{475569}

\hypersetup{
    colorlinks=true,
    linkcolor=accent,
    urlcolor=accent,
    pdftitle={DRL Lecture 1: Policy Gradient (Review) - Complete Equations}
}

\begin{document}

\begin{center}
    {\LARGE \textbf{\textcolor{accent}{DRL Lecture 1: Policy Gradient (Review)}}}\\[6pt]
    {\large Complete Mathematical Formulations and Derivations}\\[4pt]
    {\small \textcolor{muted}{Course: Deep Reinforcement Learning (Prof. Hung-yi Lee, National Taiwan University)}}
\end{center}

\vspace{6pt}
\hrule
\vspace{12pt}

\section{Reinforcement Learning Interaction and Trajectory Probability}

An actor is represented by a neural network policy $\pi$ with trainable parameters $\theta$. At time step $t$, given state (observation) $s_t$, the policy outputs a probability distribution over actions $a_t$:
\begin{equation}
    a_t \sim \pi_\theta(a_t \mid s_t) \quad \text{with action probability} \quad p_\theta(a_t \mid s_t).
\end{equation}

Over an episode of $T$ steps, the alternating interaction between the actor and the environment forms a trajectory $\tau$:
\begin{equation}
    \tau = \{s_1, a_1, s_2, a_2, \dots, s_T, a_T\}.
\end{equation}

The probability of observing trajectory $\tau$ under policy parameters $\theta$ factorizes into initial state probabilities, policy action probabilities, and environment state-transition probabilities:
\begin{align}
    p_\theta(\tau) &= p(s_1)\, p_\theta(a_1 \mid s_1)\, p(s_2 \mid s_1, a_1)\, p_\theta(a_2 \mid s_2)\, p(s_3 \mid s_2, a_2) \cdots \nonumber \\
    &= p(s_1) \prod_{t=1}^{T} p_\theta(a_t \mid s_t)\, p(s_{t+1} \mid s_t, a_t).
\end{align}

\section{Total Reward and Expected Reward Objective}

Each transition yields a scalar step reward $r_t$. The total reward accumulated over a single trajectory $\tau$ is:
\begin{equation}
    R(\tau) = \sum_{t=1}^{T} r_t.
\end{equation}

Because both the policy decisions and the environment transitions are stochastic, $R(\tau)$ is a random variable. The objective of policy gradient training is to maximize the expected total reward $\bar{R}_\theta$:
\begin{equation}
    \bar{R}_\theta = \sum_{\tau} R(\tau)\, p_\theta(\tau) = E_{\tau \sim p_\theta(\tau)}\big[R(\tau)\big].
\end{equation}

\section{Policy Gradient Derivation and Monte Carlo Estimation}

To optimize $\bar{R}_\theta$ via gradient ascent without requiring $R(\tau)$ to be differentiable, the derivation applies the log-derivative identity:
\begin{equation}
    \nabla \log f(x) = \frac{\nabla f(x)}{f(x)} \implies \nabla p_\theta(\tau) = p_\theta(\tau) \nabla \log p_\theta(\tau).
\end{equation}

Taking the gradient of the expected reward with respect to $\theta$ gives:
\begin{align}
    \nabla \bar{R}_\theta &= \sum_{\tau} R(\tau) \nabla p_\theta(\tau) \nonumber \\
    &= \sum_{\tau} R(\tau)\, p_\theta(\tau) \frac{\nabla p_\theta(\tau)}{p_\theta(\tau)} \nonumber \\
    &= \sum_{\tau} R(\tau)\, p_\theta(\tau) \nabla \log p_\theta(\tau) \nonumber \\
    &= E_{\tau \sim p_\theta(\tau)}\big[R(\tau) \nabla \log p_\theta(\tau)\big].
\end{align}

Taking the logarithm of $p_\theta(\tau)$ converts the product over time steps into a sum:
\begin{equation}
    \log p_\theta(\tau) = \log p(s_1) + \sum_{t=1}^{T} \log p_\theta(a_t \mid s_t) + \sum_{t=1}^{T} \log p(s_{t+1} \mid s_t, a_t).
\end{equation}

Because the initial state distribution $p(s_1)$ and the environment dynamics $p(s_{t+1} \mid s_t, a_t)$ do not depend on $\theta$, their gradients with respect to $\theta$ vanish:
\begin{equation}
    \nabla \log p_\theta(\tau) = \sum_{t=1}^{T} \nabla \log p_\theta(a_t \mid s_t).
\end{equation}

Approximating the expectation empirically by sampling $N$ trajectories $\{\tau^1, \tau^2, \dots, \tau^N\}$ using the current policy $\pi_\theta$, where trajectory $\tau^n$ has length $T_n$, yields the Monte Carlo policy gradient estimator:
\begin{align}
    \nabla \bar{R}_\theta &\approx \frac{1}{N} \sum_{n=1}^{N} R(\tau^n) \nabla \log p_\theta(\tau^n) \nonumber \\
    &= \frac{1}{N} \sum_{n=1}^{N} \sum_{t=1}^{T_n} R(\tau^n) \nabla \log p_\theta(a_t^n \mid s_t^n).
\end{align}

Using learning rate $\eta$, the policy parameters are updated by gradient ascent:
\begin{equation}
    \theta \leftarrow \theta + \eta \nabla \bar{R}_\theta.
\end{equation}

\section{Implementation as Weighted Multi-Class Classification}

In automatic differentiation frameworks, policy gradient training mirrors standard multi-class classification where sampled actions $a_t^n$ serve as target labels for inputs $s_t^n$.

For unweighted maximum log-likelihood classification, the objective and its gradient are:
\begin{equation}
    \mathcal{L}_{\text{MLE}}(\theta) = \frac{1}{N} \sum_{n=1}^{N} \sum_{t=1}^{T_n} \log p_\theta(a_t^n \mid s_t^n), \qquad
    \nabla \mathcal{L}_{\text{MLE}}(\theta) = \frac{1}{N} \sum_{n=1}^{N} \sum_{t=1}^{T_n} \nabla \log p_\theta(a_t^n \mid s_t^n).
\end{equation}

Multiplying each state-action log-probability term by its trajectory reward $R(\tau^n)$ produces the policy gradient surrogate objective and matching gradient:
\begin{equation}
    J(\theta) = \frac{1}{N} \sum_{n=1}^{N} \sum_{t=1}^{T_n} R(\tau^n) \log p_\theta(a_t^n \mid s_t^n), \qquad
    \nabla J(\theta) = \frac{1}{N} \sum_{n=1}^{N} \sum_{t=1}^{T_n} R(\tau^n) \nabla \log p_\theta(a_t^n \mid s_t^n).
\end{equation}

\section{Variance Reduction and Credit Assignment Improvements}

\subsection*{Tip 1: Subtracting a Baseline}

When all episode rewards $R(\tau^n)$ are positive, finite sampling reduces the probability of actions that happened not to be sampled. Subtracting a scalar baseline $b$ centers the reward weights so only better-than-average trajectories receive positive reinforcement:
\begin{equation}
    \nabla \bar{R}_\theta \approx \frac{1}{N} \sum_{n=1}^{N} \sum_{t=1}^{T_n} \big(R(\tau^n) - b\big) \nabla \log p_\theta(a_t^n \mid s_t^n),
\end{equation}
where the baseline is commonly estimated using the expected episode return:
\begin{equation}
    b \approx E\big[R(\tau)\big].
\end{equation}

\subsection*{Tip 2: Assigning Suitable Credit (Reward-to-Go and Discounting)}

Weighting action $a_t^n$ by the full episode reward $R(\tau^n)$ erroneously attributes rewards obtained prior to time step $t$ to action $a_t^n$. Replacing $R(\tau^n)$ with the reward-to-go from step $t$ to $T_n$ gives:
\begin{equation}
    \nabla \bar{R}_\theta \approx \frac{1}{N} \sum_{n=1}^{N} \sum_{t=1}^{T_n} \left( \sum_{t'=t}^{T_n} r_{t'}^n - b \right) \nabla \log p_\theta(a_t^n \mid s_t^n).
\end{equation}

To introduce discounting consistently, we now switch from the undiscounted objective $\bar{R}_\theta$ to the discounted episodic objective
\begin{equation}
    J_\gamma(\theta) = E_{\tau \sim p_\theta(\tau)}\left[\sum_{t=1}^{T} \gamma^{t-1} r_t\right], \qquad 0 < \gamma \le 1.
\end{equation}
Define the discounted reward-to-go from time $t$ as
\begin{equation}
    G_t^\gamma = \sum_{k=t}^{T} \gamma^{k-t} r_k.
\end{equation}
For ordinary on-policy trajectory sampling, the exact policy gradient of this objective is
\begin{equation}
    \nabla_\theta J_\gamma(\theta) = E_{\tau \sim p_\theta(\tau)}\left[\sum_{t=1}^{T} \gamma^{t-1} G_t^\gamma \, \nabla_\theta \log p_\theta(a_t \mid s_t)\right].
\end{equation}
The outer factor $\gamma^{t-1}$ discounts from the start of the episode, while $G_t^\gamma$ discounts relative to time $t$: their product assigns each future reward $r_k$ the objective's weight $\gamma^{k-1}$.

A fixed, action-independent baseline $b_t(s_t)$ can be subtracted without changing the expectation. With $N$ on-policy trajectories, this gives the Monte Carlo estimator
\begin{equation}
    \nabla_\theta J_\gamma(\theta) \approx \frac{1}{N} \sum_{n=1}^{N} \sum_{t=1}^{T_n} \gamma^{t-1} \big(G_t^{\gamma,n} - b_t(s_t^n)\big) \, \nabla_\theta \log p_\theta(a_t^n \mid s_t^n).
\end{equation}
Here $G_t^{\gamma,n}$ is the discounted reward-to-go in trajectory $n$. Treat the baseline as fixed when taking the policy gradient. Setting $\gamma=1$ recovers the undiscounted reward-to-go estimator above. Keeping the undiscounted objective while inserting $\gamma<1$ only inside reward-to-go instead gives a generally biased approximation; it is not an exact reformulation of that objective.

\subsection*{Advantage Function Formulation}

For the discounted objective, define time-dependent value and advantage functions (the time index accounts for the remaining horizon in a finite episode):
\begin{align}
    V_{\gamma,t}^\theta(s) &= E_{\pi_\theta}\big[G_t^\gamma \mid s_t=s\big], \\
    Q_{\gamma,t}^\theta(s,a) &= E_{\pi_\theta}\big[G_t^\gamma \mid s_t=s, a_t=a\big], \\
    A_{\gamma,t}^\theta(s,a) &= Q_{\gamma,t}^\theta(s,a) - V_{\gamma,t}^\theta(s).
\end{align}
The advantage measures how much better action $a$ is than the policy's average action at state $s$ and time $t$. The corresponding exact gradient retains the outer discount factor:
\begin{equation}
    \nabla_\theta J_\gamma(\theta) = E_{\tau \sim p_\theta(\tau)}\left[\sum_{t=1}^{T} \gamma^{t-1} A_{\gamma,t}^\theta(s_t,a_t) \, \nabla_\theta \log p_\theta(a_t \mid s_t)\right].
\end{equation}
In practice, use estimated advantages, for example $\hat{A}_{\gamma,t}^n = G_t^{\gamma,n} - \hat{V}_{\gamma,t}(s_t^n)$ with a fixed critic baseline, and form
\begin{equation}
    \nabla_\theta J_\gamma(\theta) \approx \frac{1}{N} \sum_{n=1}^{N} \sum_{t=1}^{T_n} \gamma^{t-1} \hat{A}_{\gamma,t}^n \, \nabla_\theta \log p_\theta(a_t^n \mid s_t^n).
\end{equation}
The advantage estimates are held fixed in the actor update; the policy gradient does not differentiate through them.

\section{Summary of Notation}

\begin{table}[h!]
\centering
\begin{tabular}{@{}ll@{}}
\toprule
\textbf{Symbol} & \textbf{Definition} \\
\midrule
$\theta$ & Parameters of the policy neural network \\
$\pi_\theta$ & Stochastic actor policy parameterized by $\theta$ \\
$s_t, a_t, r_t$ & State (observation), action, and immediate reward at time step $t$ \\
$\tau$ & Trajectory $\{s_1, a_1, s_2, a_2, \dots, s_T, a_T\}$ representing one episode \\
$T, T_n$ & Horizon length of an episode, and length of the $n$-th sampled episode \\
$N$ & Number of sampled trajectories per gradient update batch \\
$p_\theta(\tau)$ & Probability of sampling trajectory $\tau$ under policy $\pi_\theta$ \\
$R(\tau), \bar{R}_\theta$ & Total reward of trajectory $\tau$ and expected total reward under $\pi_\theta$ \\
$\eta$ & Gradient ascent learning rate \\
$b, b_t(s)$ & Scalar or time- and state-dependent action-independent baseline \\
$\gamma$ & Discount factor ($0 < \gamma \le 1$) in the discounted episodic objective \\
$J_\gamma(\theta)$ & Expected discounted return measured from the start of the episode \\
$G_t^\gamma, G_t^{\gamma,n}$ & Discounted reward-to-go from time $t$, and its realization in trajectory $n$ \\
$V_{\gamma,t}^\theta(s), Q_{\gamma,t}^\theta(s,a)$ & Discounted state and action values at time $t$ under $\pi_\theta$ \\
$A_{\gamma,t}^\theta(s,a)$ & Discounted advantage $Q_{\gamma,t}^\theta(s,a)-V_{\gamma,t}^\theta(s)$ \\
$\hat{A}_{\gamma,t}^n$ & Estimated advantage in trajectory $n$, held fixed during the actor update \\
\bottomrule
\end{tabular}
\end{table}

\end{document}